\documentclass[11pt]{article}
\usepackage{amsmath,amsfonts}
\usepackage{graphicx}
%\usepackage{xcolor}
%\definecolor{gray}{rgb}{.75,.75,.75}
%\usepackage[landscape]{geometry}
%\usepackage{hyperref}
%\usepackage[bookmarks=true, pdfborder={0 0 .5}, urlbordercolor=gray]{hyperref}

\renewcommand{\thefootnote}{\fnsymbol{footnote}}


\addtolength{\topmargin}{-.9in}
\addtolength{\textheight}{1.8in}
\addtolength{\oddsidemargin}{-.4in}
\addtolength{\textwidth}{.8in}
\pagestyle{empty}

\newcommand{\ds}{\displaystyle}
\newcommand{\ts}{\textstyle}

\newcommand{\Z}{\mathbb{Z}}
\newcommand{\R}{\mathbb{R}}
\newcommand{\C}{\mathbb{C}}
\newcommand{\EE}{\mathcal{E}}
\newcommand{\tZ}{\tilde{\Z}}

\newcommand{\Sym}{\mathrm{Sym}}

\renewcommand{\circle}[1]{{\bigcirc\hspace{-.75em}#1}}

\begin{document}

\footnote[0]{Notes by Peter Magyar \ {\tt magyar@math.msu.edu}}
\vspace{-1em}

\centerline{\bf Math 133\hspace{1.5em} \hfill{\Large\bf
Partial Fractions}\hfill Stewart \S7.4}
\vspace{1em}

\noindent 
{\bf Integrating basic rational functions.}
For a function $f(x)$, we have examined several algebraic methods\footnote{
Substitution \S4.5, Integration by Parts \S7.1, Products of Trig Functions \S7.2, Reverse Trig Substitution \S7.3
}
for finding its indefinite integral (antiderivative) $F(x)=\int f(x)\,dx$, 
which allows us to compute definite integrals $\int_a^b f(x)\,dx=F(b)-F(a)$ by the Second Fundamental Theorem. 

In this section, we will learn a special technique to integrate any {\it rational function},
meaning a quotient of two polynomials:
$$
f(x) \ =\ \frac{g(x)}{h(x)}
\ =\ 
\frac{a_mx^m + a_{m-1}x^{m-1} + \cdots + a_1x + a_0}
{b_nx^n + b_{n-1}x^{n-1} + \cdots + b_1x + b_0}\,,
$$
where $a_i, b_j$ are constant coefficients. 
We call the largest powers $m$ and $n$ the {\it degrees} of the polynomials $g(x)$ and $h(x)$, assuming that the highest coefficients $a_m,b_n\neq 0$.

We have several basic rational functions whose integrals we already know:
\begin{itemize}

\item[(i)] 
$\ds\int\! a_mx^m + \cdots + a_1x + a_0\ dx \ =\  \tfrac{a_m}{m+1}x^{m+1} + \cdots+  \tfrac{a_1}{2}x^2 + a_0x + C$. 

\item[(ii)] $\ds\int\! \frac{1}{x-a}\  dx \ = \  \ln\!|x{-}a|+C$.

\item[(iii)] $\ds\int\! \frac{1}{(x-a)^n}\  dx \ = \  -\frac{1}{(n{-}1)(x-a)^{n-1}}$ \ for $n\geq 2$.

\item[(iv)] $\ds\int\!\frac{x}{x^2 + a}\ dx 
\ =\ 
\tfrac 12\,\int\!\frac{1}{x^2 +a}\cdot 2x\, dx 
\ =\ 
\tfrac 12\ln\!|x^2{+}a|+C$

\item[(v)] $\ds\int\!\frac{1}{x^2 + a}\ dx 
\ =\ 
\tfrac{1}{\sqrt a}\,\int\! \frac{1}{(\frac x{\sqrt a})^2 + 1}\cdot \tfrac{1}{\sqrt a}\,dx 
\ =\ 
\tfrac 1{\sqrt a}\arctan(\tfrac x{\sqrt a})+C$,\, for $a>0$.

\item[(vi)] $\ds\int\!\frac{1}{(x^2 + 1)^2}\ dx$. \
Letting  \ $x = \tan(\theta)$, \ $x^2{+}1 =\sec^2(\theta)$, \ $dx=\sec^2(\theta)\,d\theta$:
$$\hspace{-5em}
\int\!\frac{1}{(x^2 + 1)^2}\ dx 
\ = \ 
\int\! \frac{1}{\sec^4(\theta)}\,\sec^2(\theta)\,d\theta
\ = \ 
\int\! \cos^2(\theta)\,d\theta
$$
$$
\ =\
\tfrac12( \theta + \sin(\theta)\cos(\theta)) + C
\ =\
\tfrac12\!\left(\!\arctan(x) + \frac{x}{x^2+1}\!\right) + C\,.
$$
We used: $\int\! \cos^2(\theta)\,d\theta = \int\! \frac12\,{+}\,\frac12\cos(2\theta)\,d\theta
=\frac12\theta \,{+}\, \frac14\sin(2\theta)=\frac12\theta \,{+}\,\frac12 \sin(\theta)\cos(\theta)$,
and (as in \S6.6) 
$\sin(\theta) = \frac{x}{\sqrt{x^2+1}}$, \ $\cos(\theta) = \frac{1}{\sqrt{x^2+1}}$ .
\end{itemize}
\vspace{.5em}

\noindent
{\bf Quadratic denominator.} With the above basic integrals,
we can integrate any rational function with numerator of degree at most 1 and
denominator of degree at most 2:
$$
\int\!\frac{px + q}{ax^2 + bx+c}\ dx\,.
$$
There are two different cases, depending on the sign of the {\it discriminant} $d=b^2-4ac$.
\\[.5em]
{\sc example:} Here is how to handle the case where $d=b^2-4ac>0$, such as:
$$
\ds\int\!\frac{x+1}{x^2 + x-2}\ dx, 
$$
where $d = 1^2-4(1)(-2)=9$. By the Quadratic Formula, 
the denominator has two real roots $x=\frac{-b\pm\sqrt{b^2-4ac}}{2a}=1,-2$,
which are the vertical asymptotes of our function: 
\\[.5em]
\centerline{
\includegraphics[width=3in]
{7-4-Asymptotes.png} 
}
\mbox{}
\\[-2em]

\noindent
We split our function into a sum of simple parts, 
each having just one vertical asymptote:
$$
\frac{x+1}{x^2 + x-2}
\ =\ 
\frac{x+1}{(x{-}1)(x{+}2)}
\ =\ 
\frac{A}{x{-}1} + \frac{B}{x{+}2}\,.
$$
This is called the {\it partial fraction expansion} of our rational function.
For any constants $A,B$, the graph of the right-hand function 
will have the same asymptotes as our original function, 
but we can actually find constants which make the
two exactly equal. 
Clearing denominators, we want $A,B$ such that:
$$
x+1 \ =\ 
A(x{+}2) + B(x{-}1) \quad \text{for all $x$.}
$$
Setting $x = 1$ gives $1+1 = A(1{+}2) + B(0)$, so $A=\frac23$;
and setting $x=-2$ gives $-2+1 = A(0) + B(-2{-}1)$, so $B=\frac13$. 
Now we can use the basic integral (ii) above:
$$
\int\!\frac{x+1}{x^2 + x-2}\,dx
\ =\ 
\int\!\frac{\frac23}{x{-}1} + \frac{\frac13}{x{+}2}\ dx
\ =\ 
\tfrac23\ln\!|x{-}1| + \tfrac13\ln\!|x{+}2| + C\,.
$$
\\[.5em]
{\sc example:} The other case is when $d = b^2-4ac<0$, such as:
$$
\int\!\frac{x+1}{x^2 + x + 1}\ dx,
$$
for which $d=1^2 - 4(1)(1) = -3$. In this case, the denominator has no real-number zeroes: $x^2+x+1> 0$, and it cannot be factored; it is an {\it irreducible} polynomial. The graph of $\frac{x+1}{x^2 + x + 1}$ has no vertical asymptotes.
Our strategy is to reduce its integral to the basic integrals (iv) and (v) above.

The first step is to complete the square in the denominator and force it into the form $(x{+}p)^2 + q$,
the same process that produces the Quadratic Formula:
$$
x^2 + x+1 
\ =\ 
x^2 + 2(\tfrac12)x 
+ (\tfrac12)^2 - (\tfrac12)^2 + 1
\ =\ 
(x{+}\tfrac12)^2 + \tfrac34\,.
$$
Thus, letting $u=x+\frac12$,\, $du=dx$:
$$\begin{array}{rcl}
\ds\int\!\frac{x+1}{x^2 + x+1}\ dx 
& = &
\ds\int\! \frac{(x{+}\tfrac12) -\tfrac12+1}{(x{+}\tfrac12)^2+\tfrac34}\ dx
\\[1.5em]
& = &
\ds\int\! \frac{u}{u^2+\tfrac34}\ dx
+
\tfrac12\int\! \frac{1}{u^2+\tfrac34}\ dx
\\[1.5em]
& = & 
\tfrac12\ln\!|u^2{+} \tfrac34| 
\,+\, 
\tfrac{1}{ 2\sqrt{3/4} } \arctan( \tfrac{u}{\sqrt{3/4}} ) + C
\\[1em]
& = &
\tfrac12\ln\!\left|x^2{+}x{+}1\right| 
\,+\, 
\tfrac{1}{ \sqrt{3} } \arctan( \tfrac{2x{+}1}{\sqrt3} ) + C.
\end{array}$$
In fact, the two terms in our answer correspond to splitting the original function (blue)
into a graph with reflection symmetry across the line $x=-\frac12$ (green), 
and a graph with $180^\circ$ rotation symmetry around the point $(-\frac12, 0)$ (red):
\\[1em]
\centerline{
\includegraphics[width=4in]
{7-4-Bell.png} 
}
\vspace{0em}

%
%$$\begin{array}{rcl}
%\ds\int\!\frac{px + q}{ax^2 + bx+c}\ dx 
%& = &
%\ds\int \frac{p\!\left(x+\tfrac{b}{2a}\right) -\tfrac{pb}{2a} + q}{a\!\left(\left(x+\tfrac{b}{2a}\right)^{\!2}- \tfrac{d}{4a^2}\right)}\ dx
%\\[2em]
%&& \quad\text{Let:}\quad
% u = x+\tfrac{b}{2a}  \quad \text{and} \quad h=-d>0.
%\\[1em]
%& = &
%\ds\tfrac pa\int\! \frac{u} {u^2 + \tfrac{h}{4a^2}}\ du
%\,+\, \left(\!\tfrac{2qa-pb}{2a^2}\!\right)\!\int\! \frac{ 1} {u^2 + \tfrac{h}{4a^2}}\ du
%\\[2em]
%& = &
%\ds\tfrac p{2a}\ln\left|\left(x{+}\tfrac{b}{2a}\right)^{\!2}+\tfrac{h}{4a^2}\right|
%\,+\, \left(\!\tfrac{2qa-pb}{2a^2\sqrt{h/4a^2}}\!\right)\!\arctan\!\left(\tfrac{x}%{\sqrt{h/4a^2}}\right)
%\\[2em]
%& = &
%\ds\tfrac p{2a}\ln\left|\left(x{+}\tfrac{b}{2a}\right)^{\!2}+\tfrac{4ac-b^2}{4a^2}\right|
%\,+\, \left(\!\tfrac{2qa-pb}{a\sqrt{4ac-b^2}}\!\right)\!\arctan\!\left(\tfrac{2ax}{\sqrt{4ac-b^2}}\right)\!.
%\end{array}$$
%That is quite a formula, but it is completely general for $d<0$.

\noindent
{\sc example:} One more case: if the numerator has degree greater than or equal to the denominator, for example:
$$
\int\!\frac{x^4+2x+3}{x^2+x-2}\ dx\,.
$$
Then $y=0$ is no longer a horizontal asymptote. 
Instead, the behavior of the function as $x\to\pm\infty$ is
controlled by a polynomial curve obtained by polynomial long division.
$$\begin{array}{c@{\ }c@{\!}r@{\,}c@{\,}r@{\,}c@{\,}r@{\,}c@{\,}r@{\,}c@{\,}rc}
&&&&&&x^2&-&x\,&+&3 & \text{ rem } -3x + 9
\\[-.75em]
&&\underline{\hspace{1.9in}} \hspace{-1.60in}
\\[-.15em]
x^2+x-2 & \left)\vphantom{\ds\sum}\right. \!\!\!&\ \ \ x^4 & \ \ \ &\ \ \ &\ \ \ &\ \ \ &\ + \ &2x & + & 3
\\
&&\!-(x^4 & + & x^3 & - & 2x^2 )\!\!
\\[-.75em]
&&\underline{\hspace{1.15in}} \hspace{-.95in}
\\[.4em]
&&&-&x^3 & + &2x^2 & +  &2x & + & 3 
\\[.3em]
&&& - & \!\!(x^3 &-& x\ & +& 2x)\!\!
\\[-.75em]
&&&&\underline{\hspace{1.15in}} \hspace{-.95in}
\\[.4em]
&&&&&&3x^2&&&+&3
\\
&&&&&-& \!\!(3x^2&+&3x&-&6)\!\!
\\[-.75em]
&&&&&&\underline{\hspace{1.05in}} \hspace{-.75in}
\\[.4em]
&&&&&&&-&3x&+&9
\end{array}$$
Thus $x^4+2x+3 = (x^2{-}x{+}3){\cdot}(x^2{+}x{-}2) + (-3x{+}9)$, and:
$$
\frac{x^4+2x+3}{x^2+x-2} 
\ \ =\ \
(x^2{-}x{+}3) + \frac{-3x+9}{x^2+x+1}
\ \ =\ \
(x^2{-}x{+}3) + \frac{2}{x-1} - \frac{5}{x+2}
$$
The last equality is a partial fraction expansion similar to our first example above.
Now:
$$
\int\!\frac{x^4+2x+3}{x^2+x-2} \ dx
\ =\ 
\tfrac13 x^3{-}\tfrac12 x^2{+}3x + 2\ln\!|x{-}1| -  5\ln\!|x{+}2| + C\,.
$$
\\[0em]
{\bf General case.} The above techniques suffice to integrate any rational function
$\int \frac{g(x)}{h(x)}\ dx$, provided we can factor the denominator $h(x)$.
First, we perform a partial fraction decomposition of $f(x)$ into a sum of terms of the following forms:
\begin{itemize}

\item A polynomial $q(x)$, which is the quotient in the long division $g(x)\div h(x) = q(x)$ with remainder $r(x)$ of smaller degree than $h(x)$,
so that $\frac{g(x)}{h(x)} = q(x) + \frac{r(x)}{h(x)}$. 

\item For each linear factor $x-c$ of the denominator $h(x)$, suppose $(x-a)^n$ is the highest 
power which divides $h(x)$. Then we add a sum of $n$ terms:
$$\frac{A_1}{x-a}+\frac{A_2}{(x-a)^2}+\cdots+\frac{A_n}{(x-a)^n}\,.$$

\item For each irreducible quadratic factor $ax^2 + bx+c$ of $h(x)$, 
suppose $(ax^2 + bx+c)^n$ is the highest power which divides $h(x)$. 
Then we add a sum of $n$ terms:
$$\frac{B_1x + C_1}{ax^2 + bx+c}+\frac{B_2x + C_2}{(ax^2 + bx+c)^2}
+\cdots+\frac{B_nx + C_n}{(ax^2 + bx+c)^n}\,.
$$

\end{itemize}
%
Leaving aside the polynomial term $q(x)$, we set $r(x)/h(x)$ equal to the sum of all the other terms above,
then we clear the denominators and solve for all the unknown constants in the numerators as 
we did for $A,B$ in our first example above. 
(The {\it Partial Fractions Theorem} states that a unique solution always exists.)
Once this is done, we can integrate using (i)--(vi) and the above examples (see also the integrals of higher powers below).
\\[.5em]
{\sc example:} 
We find the partial fraction expansion of: 
$$
 f(x) \ = \ \frac1{x^2(x^2+1)^2}
\ =\
\frac{A_1}{x} + \frac{A_2}{x^2} + \frac{B_1x{+}C_1}{x^2+1}+ \frac{B_2x{+}C_2}{(x^2+1)^2}\,.
$$
Since the numerator has degree less than the denominator, there is no polynomial term $q(x)$.
We need to find the six constants $A_1, A_2, B_1, B_2, C_1, C_2$ which make the above equation valid.
Clearing denominators gives:
$$\begin{array}{rcl}
1 &=& A_1x(x^2{+}1)^2 + A_2(x^2{+}1)^2 + (B_1x{+}C_1)x^2(x^2{+}1) + (B_2x{+}C_2)x^2 
\\
&=& 
(A_1{+}B_1)x^5+ (A_2{+}C_1)x^4 + (2 A_1{+}B_1{+}B_2)x^3+(2 A_2{+}C_1{+}C_2)x^2+A_1 x+ A_2
\end{array}$$
Since this is an equality of polynomial functions, the coefficients of $x^k$ on the right 
must equal the coefficients of $1 = 0x^5 + \cdots +0x +1$ on the left:
$$\begin{array}{c}
A_1+B_1 = 0
\\
 A_2+C_1=0
\\
 2 A_1+B_1+B_2 = 0
\\
 2 A_2+C_1+C_2 = 0
\\
A_1 = 0, \quad A_2=1.
\end{array}$$
We solve this as:
$$
A_1 = 0, \quad A_2 = 1, \quad B_1 = -A_1 = 0, \quad C_1 = -A_2 = -1, 
$$
$$
B_2 = -2A_1 - B_1 = 0, \quad 
C_2 = -2A_2 - C_1 =  -1\,.
$$
Hence, according to (i)--(vi):
$$\begin{array}{rcl}
\ds\int\! \frac1{x^2(x^2+1)^2}\ dx
&=&
\ds\int\!\frac{1}{x^2} \, {-}\,  \frac{1}{x^2+1}\, {-}\,  \frac{1}{(x^2+1)^2}\ dx
\\[1em]
&=&
\ds-\frac1x - \arctan(x) - \frac12\!\left(\arctan(x) + \frac{x}{x^2+1}\right).
\end{array}$$

%$$\hspace{-.3in}
%f(x) \ =\  q(x) + 
%\left(\frac{g_{1,1}(x)}{h_1(x)} +
%\frac{g_{1,2}(x)}{h_1(x)^2} + \cdots +
%\frac{g_{1,n_1}(x)}{h_1(x)^{n_1}}\right)
%$$
%$$
%\hspace{1.2in} +\ \cdots\  + 
%\left(\frac{g_{k,1}(x)}{h_k(x)} +
%\frac{g_{k,2}(x)}{h_k(x)^2} + \cdots +
%\frac{g_{k,n_k}(x)}{h_k(x)^{n_k}}
%\right)\,,
%$$
%where $q(x)$ is the quotient polynomial of $g(x)/h(x)$ and each
%$g_{i,j}(x)$ is either a constant (if $h_i(x)$ is linear)
%or a  linear polynomial (of $h_i(x)$ is quadratic).
%Note that the total number of terms is $n_1+\cdots+n_k=n$, the degree of the denominator $h(x)$,
%and each term has an unknown constant or quadratic  $g_{i,j}(x)$ which must be solved for like $A,B$ in our previous example.
%Once this is done, we can integrate using (i)--(vi) and the above examples.\footnote{
%For $\int\!\frac{1}{(x^2 + 1)^j}\,dx$ we need even more elaborate contortions.
%The truth is this only becomes manageable if we use imaginary number factorizations like $x^2+1 =(x{-}\sqrt{-1})(x{+}\sqrt{-1})$, avoiding quadratic denominators entirely.
%}

\pagebreak

\noindent
{\bf Trig integrals again.}  In \S7.2--7.3, we reduced trig integrals by substitution 
to rational function integrals, which we can now find by partial fractions.
For example:
$$
\int\!\sec(x)\,dx 
\ =\
\int\! \frac{1}{\cos^2(x)}\cdot \cos(x)\,dx 
\ =\
\int\! \frac{1}{1-\sin^2(x)}\cdot\cos(x)\,dx
$$
$$
\left\{\begin{array}{r@{\ }c@{\ }l} u &=& \sin(x) \\ du &=& \cos(x)\,dx\end{array}\right\}
\quad
\ =\
\int\! \frac1{1-u^2}\ du
\ =\
\int\!\frac{\tfrac12}{1+u} + \frac{\tfrac12}{1-u}\ du
$$
$$
\ =\
\tfrac12\ln(1{+}u)-\tfrac12\ln(1{-}u)
\ =\
\ln  \sqrt{\frac{1{+}u}{1{-}u}}
\ =\
\ln  \sqrt{\frac{1{+}\sin(x)}{1{-}\sin(x)}}\ .
$$
\\[.5em]
{\sc challenge:} Show by identities that this is equal to our previous answer 
$\int\sec(x)\,dx=\ln\!|\!\tan(x)+\sec(x)|$ given in \S7.2 .
Also: try this method on $\int\!\sec^3(x)\,dx=\int\!\! \frac{1}{(1-u^2)^2}\,du$.
\\[.5em]
{\sc challenge:}  Integrate $\int\! \sqrt{\tan(x)}\,dx$. 
The substitution $u=\tan(x)$, $x=\arctan(u)$, $dx = \frac1{u^2+1}\,du$ 
gives $\int\! \frac{\sqrt u}{u^2+1}\,du$. Then $z=\sqrt u$ reduces this to a rational function.
For partial fractions, factor $z^4+1=(z^2{+}1)^2-2z^2=(z^2+1+\sqrt2\, z)(z^2+1-\sqrt2\,z)$.
Solve:
$$%\textstyle
\frac{z^2}{z^4+1} \ =\ \frac{Az+B}{z^2+\sqrt2\, z+1} + \frac{Cz+D}{z^2-\sqrt2\, z+1}\, .
$$
%
{\bf Higher quadratic powers.} For powers of irreducible quadratics in the denominator:
  $$
  \hspace{-.5in} (*) \qquad
 \int\!\frac{x}{(x^2 + 1)^n}\,dx \ =\ -\frac{1}{2(n{-}1)(1+x^2)^{n-1}},
 $$
using the substitution $u=x^2+1$. 
The other basic rational integral
$$
I_n=\int\!\frac{1}{(x^2 + 1)^n}\,dx
$$ 
can be evaluated {\it recursively} using Integration by Parts and the previous integral:
$$\hspace{-.1in} 
\begin{array}{r@{\ }ll}
I_n =&   \int\!\frac{1}{(x^2 + 1)^n}\,dx 
\ =\ 
\int \frac{x^2 + 1}{(x^2 + 1)^n}\,dx 
- \int \frac{x^2}{(x^2 + 1)^n}\,dx 
& \text{since} \ 1 = (x^2{+}1) -x^2
\\[.8em]
=&
\int \frac{1}{(x^2 + 1)^{n-1}}\,dx 
- \int x\cdot\frac{x}{(x^2 + 1)^n}\,dx 
& \text{setting up Int By Parts}
\\[.8em]
=&
\int \frac{1}{(x^2 + 1)^{n-1}}\,dx 
 + x\cdot \frac{1}{2(n-1)(x^2 + 1)^{n-1}} 
 -  \int 1\cdot \frac{1}{2(n-1)(x^2 + 1)^{n-1}}\,dx
& \text{by IBP and  integral $(*)$}
\\[.8em]
=&
 I_{n-1} + \frac{x}{2(n-1)(1+x^2)^{n-1}} - \frac{1}{2(n-1)}I_{n-1}
& \text{by definition of $I_{n-1}$.}
\end{array}$$
Simplifying, we get the recursive formula:
$$
I_n \ \ =\ \ \frac{2n{-}3}{2(n{-}1)}I_{n-1} \,+\, \frac{x}{2(n{-}1)(x^2 + 1)^{n-1}}.
$$
{\sc example:} For $n=3$, we already know $I_2 = \int\!\frac{1}{(x^2 + 1)^2}\,dx
= \frac12(\arctan(x) +\frac{x}{x^2+1})$, so:
$$\begin{array}{rcl}
I_3 &=& \int\!\frac{1}{(x^2 + 1)^3}\,dx 
\ =\  \frac{2(3)-3}{2(3-1)}I_{2} + \frac{x}{2(3-1)(x^2 + 1)^{2}}
\\[.8em]
&=&
  \frac34 \frac12(\arctan(x) +\frac{x}{x^2+1})+ \frac14\frac{x}{(x^2 + 1)^{2}}
\\[.8em]
&=&
  \frac38\arctan(x) +\frac{x (3 x^2 + 5)}{8 (x^2 + 1)^2}\,.
\end{array}$$


\end{document}


We need even more elaborate contortions.
The truth is this only becomes manageable if we use imaginary number factorizations like $x^2+1 =(x{-}\sqrt{-1})(x{+}\sqrt{-1})$, avoiding quadratic denominators entirely.


%%%%%%%%%%%  Picture  %%%%%%%%%%
\\[0em]
\centerline{
%\includegraphics[width=2in]
{1-8-Discont-iv.png}
\qquad \qquad 
%\includegraphics[width=2in]
{1-8-Discont-v.png} 
}
\vspace{0em}

\noindent
%%%%%%%%%%%%%%%%%%%%%%%%%

%%%%%%%%%%%  Table  %%%%%%%%%%
In fact, we get the following derivatives:
$$\begin{array}{|c||c|c|c|c|c|c|}
\hline &&&&&& \\[-.9em] 
f(x) & \sin(x) & \cos(x) & \tan(x)  & \sec(x) & \csc(x) & \cot(x)
\\[.2em] \hline &&&&&& \\[-.9em]
f'(x) &\cos(x) & -\sin(x) & \sec^2(x) & \tan(x)\sec(x) & -\cot(x)\csc(x) & -\csc^2(x)
\\[.1em] \hline
\end{array}$$
\\[1em]
%%%%%%%%%%%%%%%%%%%%%%%%%%

























